% Shared body of questions.tex / questions-refs.tex.
% Placeholder excerpt (from the 2025 зачёт list) until the syllabus is final.
\examtitle[вопросов на зачёт]{Вопросы на зачёт}{декабрь 2026}

\begin{literature}
  \item Д.~Э.~Кнут. Искусство программирования, Т2.
  \item М.~Вельшенбах. Криптография на Си и С++ в действии.
  \item Р.~Крэндалл, К.~Померанс. Простые числа: криптографические и вычислительные аспекты.
\end{literature}

\topic{Арифметика больших чисел и операции в $\Zn{n}$}
\question{Школьные алгоритмы сложения, умножения и деления с остатком больших целых чисел, оценки сложности.}{[1, Раздел 4.3], [2, Глава 4]}
\question{Бинарный алгоритм быстрого возведения в степень и его модификации.}{[2, Глава 6]}
\question{Теорема Монтгомери и алгоритм умножения чисел по модулю (без тривиальных делений).}{[3, Глава 9]}

\topic{Простые числа}
\question{Свидетели (не)простоты, вероятностный тест простоты Миллера"--~Рабина.}{[3, Глава 3]}
